\documentclass[11pt,a4paper]{article}
\usepackage[utf8]{inputenc}
\usepackage[spanish]{babel}
\usepackage[T1]{fontenc}
\usepackage{geometry}
\geometry{margin=2cm}
\usepackage{xcolor}
\usepackage{listings}
\usepackage{tcolorbox}
\tcbuselibrary{most}
\usepackage{booktabs}
\usepackage{array}
\usepackage{hyperref}
\usepackage{fancyhdr}
\usepackage{titlesec}
\usepackage{multicol}
\usepackage{amssymb}

% Colores
\definecolor{pyblue}{RGB}{53,114,165}
\definecolor{tsblue}{RGB}{0,122,204}
\definecolor{vuegreen}{RGB}{66,184,131}
\definecolor{codebg}{RGB}{245,245,245}
\definecolor{alertred}{RGB}{220,50,50}
\definecolor{tipgreen}{RGB}{39,174,96}
\definecolor{warnambel}{RGB}{230,126,34}
\definecolor{darkgray}{RGB}{50,50,50}

% Listings Python
\lstdefinestyle{python}{
  language=Python,
  backgroundcolor=\color{codebg},
  basicstyle=\ttfamily\small,
  keywordstyle=\color{pyblue}\bfseries,
  stringstyle=\color{teal},
  commentstyle=\color{gray}\itshape,
  numbers=left, numberstyle=\tiny\color{gray},
  frame=single, framerule=0.4pt, rulecolor=\color{pyblue!40},
  breaklines=true, tabsize=2,
  showstringspaces=false,
  morekeywords={async,await,def,class,from,import,as,TypeVar,Generic,Optional,Union,List,Dict}
}

% Listings TypeScript
\lstdefinestyle{typescript}{
  language=Java,
  backgroundcolor=\color{codebg},
  basicstyle=\ttfamily\small,
  keywordstyle=\color{tsblue}\bfseries,
  stringstyle=\color{teal},
  commentstyle=\color{gray}\itshape,
  numbers=left, numberstyle=\tiny\color{gray},
  frame=single, framerule=0.4pt, rulecolor=\color{tsblue!40},
  breaklines=true, tabsize=2,
  showstringspaces=false,
  morekeywords={interface,type,string,number,boolean,void,any,never,readonly,
                const,let,async,await,export,import,from,extends,implements,
                ref,reactive,computed,watch,onMounted,defineComponent,Ref,
                Promise,Record,Partial,Required,Pick,Omit,Array,Generic}
}

% Cajas
\newtcolorbox{tipbox}{colback=tipgreen!10,colframe=tipgreen,title=\textbf{Tip},fonttitle=\small\bfseries}
\newtcolorbox{alertbox}{colback=alertred!10,colframe=alertred,title=\textbf{Ojo},fonttitle=\small\bfseries}
\newtcolorbox{pybox}{colback=pyblue!8,colframe=pyblue,title=\textbf{Python},fonttitle=\small\bfseries}
\newtcolorbox{tsbox}{colback=tsblue!8,colframe=tsblue,title=\textbf{TypeScript},fonttitle=\small\bfseries}
\newtcolorbox{vuebox}{colback=vuegreen!10,colframe=vuegreen,title=\textbf{Vue.js},fonttitle=\small\bfseries}
\newtcolorbox{ejerciciobox}[1]{colback=warnambel!8,colframe=warnambel,title=\textbf{Ejercicio #1},fonttitle=\small\bfseries}
\newtcolorbox{solucionbox}[1]{colback=tipgreen!8,colframe=tipgreen!60,title=\textbf{Solución #1},fonttitle=\small\bfseries}

% Headers
\pagestyle{fancy}
\fancyhf{}
\lhead{\textcolor{tsblue}{\textbf{TypeScript para Pythonistas}}}
\rhead{\textcolor{gray}{Preparación Técnica — Toku}}
\cfoot{\thepage}

% Secciones en color
\titleformat{\section}{\large\bfseries\color{tsblue}}{}{0em}{}[\titlerule]
\titleformat{\subsection}{\normalsize\bfseries\color{darkgray}}{}{0em}{}

\begin{document}

% ─── PORTADA ─────────────────────────────────────────────────────────────────
\begin{titlepage}
  \centering
  \vspace*{2cm}
  {\Huge\bfseries\color{tsblue} TypeScript para Pythonistas\par}
  \vspace{0.5cm}
  {\large\color{gray} Preparación Técnica — Toku\par}
  \vspace{0.3cm}
  \textcolor{vuegreen}{\rule{8cm}{2pt}}
  \vspace{1cm}

  \begin{tcolorbox}[colback=tsblue!5,colframe=tsblue,width=10cm]
    \centering
    \textbf{Stack objetivo:} \texttt{TypeScript + Vue.js + FastAPI}\\[4pt]
    \textbf{Foco:} Entender TS viniendo de Python\\[4pt]
    \textbf{Nivel:} Intermedio Python → Básico-Intermedio TS
  \end{tcolorbox}

  \vspace{1.5cm}
  {\large Temas\par}
  \vspace{0.3cm}
  \begin{tabular}{ll}
    \textcolor{tsblue}{\textbullet} & Tipado estático y estructural\\
    \textcolor{tsblue}{\textbullet} & Funciones y arrow functions\\
    \textcolor{tsblue}{\textbullet} & Async/Await y Promises\\
    \textcolor{tsblue}{\textbullet} & Generics\\
    \textcolor{tsblue}{\textbullet} & Módulos e imports\\
    \textcolor{vuegreen}{\textbullet} & Vue.js Composition API\\
    \textcolor{warnambel}{\textbullet} & 5 Ejercicios con soluciones\\
  \end{tabular}

  \vfill
  {\small\color{gray} Marzo 2026 — Estudio Personal\par}
\end{titlepage}

\tableofcontents
\newpage

% ─── SECCIÓN 1: TIPADO ───────────────────────────────────────────────────────
\section{Tipado: De Python a TypeScript}

\subsection{Tipos primitivos}

\begin{tcolorbox}[colback=gray!5,colframe=gray,title=La idea clave]
Python tiene \textbf{type hints} opcionales (\texttt{int}, \texttt{str}...). TypeScript tiene tipado \textbf{estático obligatorio}
(en la práctica). Si no pones tipo, TS infiere; si infiere \texttt{any}, es una señal de alerta.
\end{tcolorbox}

\begin{pybox}
\begin{lstlisting}[style=python]
# Python
nombre: str = "Fabian"
edad: int = 25
precio: float = 9990.0
activo: bool = True
nada: None = None
cualquier_cosa: any = 42   # evitar
\end{lstlisting}
\end{pybox}

\begin{tsbox}
\begin{lstlisting}[style=typescript]
// TypeScript
const nombre: string = "Fabian";
const edad: number = 25;        // number = int y float
const precio: number = 9990.0;
const activo: boolean = true;
const nada: null = null;
const indefinido: undefined = undefined;
let cualquier_cosa: any = 42;   // evitar: apaga el tipado
\end{lstlisting}
\end{tsbox}

\begin{alertbox}
En TS \texttt{number} cubre tanto enteros como flotantes (no hay \texttt{int} y \texttt{float} separados).
\end{alertbox}

\subsection{Interfaces y Types}

Equivalente a las \texttt{dataclass} o \texttt{TypedDict} de Python:

\begin{pybox}
\begin{lstlisting}[style=python]
from dataclasses import dataclass
from typing import TypedDict, Optional

@dataclass
class Cliente:
    id: str
    nombre: str
    email: str
    telefono: Optional[str] = None

class ClienteDict(TypedDict):
    id: str
    nombre: str
    email: str
\end{lstlisting}
\end{pybox}

\begin{tsbox}
\begin{lstlisting}[style=typescript]
// Interface (preferido para objetos/formas)
interface Cliente {
  id: string;
  nombre: string;
  email: string;
  telefono?: string;    // ? = opcional (equivale a Optional[str])
  readonly createdAt: Date; // readonly = no mutable
}

// Type alias (mas flexible, puede ser union/intersection)
type Estado = "activo" | "inactivo" | "pendiente"; // union literal

type ClienteResumen = Pick<Cliente, "id" | "nombre">; // solo esos campos
type ClienteParcial = Partial<Cliente>;               // todos opcionales
\end{lstlisting}
\end{tsbox}

\begin{tipbox}
\textbf{interface vs type}: Usa \texttt{interface} para describir objetos y clases.
Usa \texttt{type} para uniones, intersecciones y alias de tipos complejos.
En la práctica son muy similares para casos simples.
\end{tipbox}

\subsection{Union e Intersection}

\begin{pybox}
\begin{lstlisting}[style=python]
from typing import Union

def procesar(valor: Union[str, int]) -> str:
    return str(valor)

# Python 3.10+
def procesar(valor: str | int) -> str:
    return str(valor)
\end{lstlisting}
\end{pybox}

\begin{tsbox}
\begin{lstlisting}[style=typescript]
// Union: puede ser uno U otro
type IdOrName = string | number;

function procesar(valor: string | number): string {
  return String(valor);
}

// Intersection: debe cumplir ambos
interface Timestamps { createdAt: Date; updatedAt: Date; }
interface ClienteBase { id: string; nombre: string; }

type ClienteCompleto = ClienteBase & Timestamps;
// ahora necesita id, nombre, createdAt Y updatedAt
\end{lstlisting}
\end{tsbox}

% ─── SECCIÓN 2: FUNCIONES ────────────────────────────────────────────────────
\section{Funciones y Arrow Functions}

\begin{pybox}
\begin{lstlisting}[style=python]
# Funcion normal
def saludar(nombre: str, saludo: str = "Hola") -> str:
    return f"{saludo}, {nombre}!"

# Lambda
cuadrado = lambda x: x ** 2

# Funcion que recibe funcion
def aplicar(fn: callable, valor: int) -> int:
    return fn(valor)
\end{lstlisting}
\end{pybox}

\begin{tsbox}
\begin{lstlisting}[style=typescript]
// Funcion normal (tipada)
function saludar(nombre: string, saludo: string = "Hola"): string {
  return `${saludo}, ${nombre}!`;  // template literals con backtick
}

// Arrow function (equivale a lambda pero mas poderosa)
const cuadrado = (x: number): number => x ** 2;

// Arrow sin llaves = return implicito (solo expresiones simples)
const doble = (x: number) => x * 2;

// Tipo de una funcion
type Transformador = (valor: number) => number;
const aplicar = (fn: Transformador, valor: number): number => fn(valor);

// Parametros rest (equivale a *args)
const sumar = (...numeros: number[]): number =>
  numeros.reduce((acc, n) => acc + n, 0);
\end{lstlisting}
\end{tsbox}

\begin{tipbox}
Las arrow functions en TS/JS tienen una diferencia crucial con \texttt{def}: el valor de
\texttt{this} es léxico (del contexto donde se define). Esto importa en Vue.js y clases.
\end{tipbox}

% ─── SECCIÓN 3: ASYNC/AWAIT ──────────────────────────────────────────────────
\section{Async/Await y Promises}

\subsection{Comparación directa}

\begin{pybox}
\begin{lstlisting}[style=python]
import asyncio
import httpx

async def obtener_cliente(id: str) -> dict:
    async with httpx.AsyncClient() as client:
        response = await client.get(f"/customers/{id}")
        return response.json()

async def main():
    cliente = await obtener_cliente("cust_123")
    print(cliente["nombre"])

asyncio.run(main())
\end{lstlisting}
\end{pybox}

\begin{tsbox}
\begin{lstlisting}[style=typescript]
// Promise<T> = equivale a Coroutine/Awaitable en Python
async function obtenerCliente(id: string): Promise<Cliente> {
  const response = await fetch(`/customers/${id}`);
  if (!response.ok) throw new Error("Cliente no encontrado");
  return response.json() as Promise<Cliente>;
}

// Uso
async function main() {
  const cliente = await obtenerCliente("cust_123");
  console.log(cliente.nombre);
}

// Manejo de errores (equivale a try/except)
async function seguro() {
  try {
    const cliente = await obtenerCliente("cust_123");
  } catch (error) {
    console.error("Error:", error);
  } finally {
    console.log("Siempre se ejecuta");
  }
}
\end{lstlisting}
\end{tsbox}

\subsection{Promises en paralelo}

\begin{pybox}
\begin{lstlisting}[style=python]
import asyncio

# Ejecutar en paralelo
clientes = await asyncio.gather(
    obtener_cliente("c1"),
    obtener_cliente("c2"),
    obtener_cliente("c3"),
)
\end{lstlisting}
\end{pybox}

\begin{tsbox}
\begin{lstlisting}[style=typescript]
// Promise.all = asyncio.gather
const [c1, c2, c3] = await Promise.all([
  obtenerCliente("c1"),
  obtenerCliente("c2"),
  obtenerCliente("c3"),
]);

// Promise.allSettled = no falla si uno falla
const resultados = await Promise.allSettled([
  obtenerCliente("c1"),
  obtenerCliente("inexistente"),  // este puede fallar
]);
resultados.forEach(r => {
  if (r.status === "fulfilled") console.log(r.value);
  else console.error(r.reason);
});
\end{lstlisting}
\end{tsbox}

% ─── SECCIÓN 4: GENERICS ─────────────────────────────────────────────────────
\section{Generics}

\begin{pybox}
\begin{lstlisting}[style=python]
from typing import TypeVar, Generic, List

T = TypeVar('T')

class Repositorio(Generic[T]):
    def __init__(self):
        self._items: List[T] = []

    def agregar(self, item: T) -> None:
        self._items.append(item)

    def obtener(self, i: int) -> T:
        return self._items[i]

repo_clientes: Repositorio[Cliente] = Repositorio()
\end{lstlisting}
\end{pybox}

\begin{tsbox}
\begin{lstlisting}[style=typescript]
// Generic con <T> (sintaxis directa, sin TypeVar)
class Repositorio<T> {
  private items: T[] = [];

  agregar(item: T): void {
    this.items.push(item);
  }

  obtener(i: number): T {
    return this.items[i];
  }
}

const repoClientes = new Repositorio<Cliente>();

// Funciones genericas
function primero<T>(lista: T[]): T | undefined {
  return lista[0];
}

// Con constraint (equivale a TypeVar con bound=)
function obtenerNombre<T extends { nombre: string }>(entidad: T): string {
  return entidad.nombre;
}

// Tipos de utilidad genericos (muy usados en APIs)
type ApiResponse<T> = {
  data: T;
  ok: boolean;
  error?: string;
};

type ClienteResponse = ApiResponse<Cliente>;
type ListaResponse = ApiResponse<Cliente[]>;
\end{lstlisting}
\end{tsbox}

% ─── SECCIÓN 5: MÓDULOS ──────────────────────────────────────────────────────
\section{Módulos e Imports}

\begin{pybox}
\begin{lstlisting}[style=python]
# cliente.py
class Cliente:
    pass

def crear_cliente(data: dict) -> Cliente:
    pass

# main.py
from cliente import Cliente, crear_cliente
import cliente as mod
from cliente import *  # evitar
\end{lstlisting}
\end{pybox}

\begin{tsbox}
\begin{lstlisting}[style=typescript]
// cliente.ts
export interface Cliente {
  id: string;
  nombre: string;
}

export function crearCliente(data: Partial<Cliente>): Cliente {
  return { id: crypto.randomUUID(), nombre: data.nombre ?? "" };
}

export default class ClienteService {}  // export default = un solo default

// main.ts
import { Cliente, crearCliente } from "./cliente";  // named exports
import ClienteService from "./cliente";              // default export
import * as clienteMod from "./cliente";             // namespace import
import type { Cliente } from "./cliente";            // solo el tipo (no JS)
\end{lstlisting}
\end{tsbox}

\begin{tipbox}
\texttt{import type} es importante en TypeScript: importa solo la definición de tipo,
sin incluir el código JavaScript en el bundle. Úsalo cuando solo necesitas el tipo
para anotaciones, no el valor en runtime.
\end{tipbox}

% ─── SECCIÓN 6: VUE.JS ───────────────────────────────────────────────────────
\section{Vue.js 3 — Composition API}

\begin{vuebox}
\begin{lstlisting}[style=typescript]
<script setup lang="ts">
// <script setup> es la forma moderna (Vue 3.2+)
// Equivale a defineComponent() pero mas limpio

import { ref, reactive, computed, watch, onMounted } from 'vue';

// ref() = variable reactiva simple (primitivos)
// Acceso con .value en JS, pero sin .value en el template
const count = ref<number>(0);
const nombre = ref<string>('');

// reactive() = objeto reactivo completo
const estado = reactive({
  cargando: false,
  error: null as string | null,
  clientes: [] as Cliente[],
});

// computed() = propiedad derivada (equivale a @property en Python)
const nombreMayusculas = computed(() => nombre.value.toUpperCase());
const hayClientes = computed(() => estado.clientes.length > 0);

// watch() = observar cambios (como un setter con efecto secundario)
watch(nombre, (nuevoValor, valorAnterior) => {
  console.log(`Cambio: ${valorAnterior} -> ${nuevoValor}`);
});

// onMounted() = se ejecuta cuando el componente monta (como useEffect[] en React)
onMounted(async () => {
  estado.cargando = true;
  try {
    estado.clientes = await obtenerClientes();
  } finally {
    estado.cargando = false;
  }
});

// Funciones normales (automaticamente disponibles en el template)
function incrementar() {
  count.value++;
}

async function cargarCliente(id: string) {
  const cliente = await obtenerCliente(id);
  estado.clientes.push(cliente);
}
</script>

<template>
  <!-- count sin .value en el template -->
  <p>{{ count }}</p>
  <button @click="incrementar">+1</button>

  <!-- v-if, v-for son equivalentes a if/for de Jinja2 -->
  <div v-if="estado.cargando">Cargando...</div>
  <ul v-else>
    <li v-for="cliente in estado.clientes" :key="cliente.id">
      {{ cliente.nombre }}
    </li>
  </ul>
</template>
\end{lstlisting}
\end{vuebox}

\subsection{Props y Emits (comunicación entre componentes)}

\begin{vuebox}
\begin{lstlisting}[style=typescript]
<script setup lang="ts">
// Props (datos que entran al componente, como argumentos)
interface Props {
  clienteId: string;
  readonly?: boolean;
}
const props = defineProps<Props>();

// Emits (eventos que salen del componente)
const emit = defineEmits<{
  'cliente-guardado': [cliente: Cliente];
  'cancelado': [];
}>();

// Uso:
emit('cliente-guardado', { id: '1', nombre: 'Juan', email: 'j@j.com' });
</script>

<template>
  <!-- :clienteId pasa prop dinamica, clienteId="literal" es string literal -->
  <ChildComponent :cliente-id="props.clienteId" @cliente-guardado="manejar" />
</template>
\end{lstlisting}
\end{vuebox}

% ─── SECCIÓN 7: TABLA DE EQUIVALENCIAS ──────────────────────────────────────
\section{Tabla de Equivalencias Python $\to$ TypeScript}

\begin{center}
\begin{tabular}{p{4cm} p{4.5cm} p{5cm}}
\toprule
\textbf{Concepto} & \textbf{Python} & \textbf{TypeScript} \\
\midrule
Tipo string & \texttt{str} & \texttt{string} \\
Tipo entero & \texttt{int} & \texttt{number} \\
Tipo flotante & \texttt{float} & \texttt{number} \\
Tipo booleano & \texttt{bool} & \texttt{boolean} \\
Sin tipo / desconocido & \texttt{Any} & \texttt{any} / \texttt{unknown} \\
Opcional & \texttt{Optional[T]} & \texttt{T | undefined} / \texttt{T?} \\
Lista & \texttt{List[T]} & \texttt{T[]} / \texttt{Array<T>} \\
Diccionario & \texttt{Dict[str, T]} & \texttt{Record<string, T>} \\
Tupla & \texttt{Tuple[int, str]} & \texttt{[number, string]} \\
Clase de datos & \texttt{@dataclass} & \texttt{interface} / \texttt{class} \\
Herencia & \texttt{class B(A)} & \texttt{class B extends A} \\
Mixin & \texttt{class C(A, B)} & \texttt{class C extends A implements B} \\
TypeVar/Generic & \texttt{TypeVar, Generic[T]} & \texttt{<T>} \\
Lambda & \texttt{lambda x: x+1} & \texttt{(x) => x+1} \\
Comprensión lista & \texttt{[x*2 for x in l]} & \texttt{l.map(x => x*2)} \\
Filtro & \texttt{[x for x in l if x>0]} & \texttt{l.filter(x => x>0)} \\
Reduce & \texttt{reduce(fn, l, 0)} & \texttt{l.reduce(fn, 0)} \\
Formato string & \texttt{f"hola \{x\}"} & \texttt{`hola \${x}`} \\
Null & \texttt{None} & \texttt{null} / \texttt{undefined} \\
Asyncio gather & \texttt{asyncio.gather()} & \texttt{Promise.all()} \\
Try/except/finally & \texttt{try/except/finally} & \texttt{try/catch/finally} \\
Decorador & \texttt{@decorator} & \texttt{@Decorator()} (limitado) \\
Módulo import & \texttt{from mod import fn} & \texttt{import \{ fn \} from './mod'} \\
\bottomrule
\end{tabular}
\end{center}

% ─── SECCIÓN 8: EJERCICIOS ───────────────────────────────────────────────────
\section{Ejercicios Prácticos}

\begin{ejerciciobox}{1 — Tipado básico}
Tienes esta función en Python. Pásala a TypeScript con tipado correcto:
\begin{lstlisting}[style=python]
def calcular_monto_con_interes(
    monto: float,
    tasa: float = 0.05,
    periodos: int = 12
) -> dict:
    total = monto * ((1 + tasa) ** periodos)
    return {"original": monto, "total": total, "interes": total - monto}
\end{lstlisting}
\end{ejerciciobox}

\begin{solucionbox}{1}
\begin{lstlisting}[style=typescript]
interface ResultadoInteres {
  original: number;
  total: number;
  interes: number;
}

function calcularMontoConInteres(
  monto: number,
  tasa: number = 0.05,
  periodos: number = 12
): ResultadoInteres {
  const total = monto * Math.pow(1 + tasa, periodos);
  return { original: monto, total, interes: total - monto };
}
\end{lstlisting}
\end{solucionbox}

\begin{ejerciciobox}{2 — Async/Await con manejo de error}
Escribe una función TypeScript \texttt{obtenerClienteSeguro(id: string)} que llame a una API,
devuelva el cliente si existe, o \texttt{null} si falla (sin lanzar error al caller).
\end{ejerciciobox}

\begin{solucionbox}{2}
\begin{lstlisting}[style=typescript]
async function obtenerClienteSeguro(id: string): Promise<Cliente | null> {
  try {
    const response = await fetch(`https://api.trytoku.com/customers/${id}`, {
      headers: { Authorization: `Bearer ${process.env.TOKU_API_KEY}` },
    });
    if (!response.ok) return null;
    return response.json() as Promise<Cliente>;
  } catch {
    return null;
  }
}
\end{lstlisting}
\end{solucionbox}

\begin{ejerciciobox}{3 — Generics con restriccion}
Crea una función genérica \texttt{paginar<T>} que reciba un arreglo de T,
un número de página y tamaño de página, y devuelva
\texttt{\{ items: T[], total: number, pagina: number, totalPaginas: number \}}.
\end{ejerciciobox}

\begin{solucionbox}{3}
\begin{lstlisting}[style=typescript]
interface Pagina<T> {
  items: T[];
  total: number;
  pagina: number;
  totalPaginas: number;
}

function paginar<T>(lista: T[], pagina: number, tamano: number = 10): Pagina<T> {
  const total = lista.length;
  const totalPaginas = Math.ceil(total / tamano);
  const inicio = (pagina - 1) * tamano;
  const items = lista.slice(inicio, inicio + tamano);
  return { items, total, pagina, totalPaginas };
}
\end{lstlisting}
\end{solucionbox}

\begin{ejerciciobox}{4 — Vue.js Composable}
Crea un composable \texttt{useClientes()} que exponga:
\texttt{clientes} (lista reactiva), \texttt{cargando} (boolean),
\texttt{cargar()} (función async) y \texttt{total} (computed).
\end{ejerciciobox}

\begin{solucionbox}{4}
\begin{lstlisting}[style=typescript]
// useClientes.ts
import { ref, computed } from 'vue';

export function useClientes() {
  const clientes = ref<Cliente[]>([]);
  const cargando = ref(false);
  const error = ref<string | null>(null);

  const total = computed(() => clientes.value.length);

  async function cargar() {
    cargando.value = true;
    error.value = null;
    try {
      const resp = await fetch('/api/customers');
      clientes.value = await resp.json();
    } catch (e) {
      error.value = 'Error al cargar clientes';
    } finally {
      cargando.value = false;
    }
  }

  return { clientes, cargando, error, total, cargar };
}
\end{lstlisting}
\end{solucionbox}

\begin{ejerciciobox}{5 — Tipos de utilidad Toku}
Dado el modelo de Toku, crea tipos TypeScript para:
\texttt{Invoice} (deuda), \texttt{InvoiceCreate} (para crear, sin id ni fechas automáticas),
\texttt{InvoiceEstado} (union literal de estados posibles).
\end{ejerciciobox}

\begin{solucionbox}{5}
\begin{lstlisting}[style=typescript]
type InvoiceEstado = "pending" | "paid" | "failed" | "void" | "partially_paid";

interface Invoice {
  id: string;
  customerId: string;
  productId: string;
  externalId?: string;
  amount: number;           // en centavos o moneda base
  currency: "CLP" | "MXN" | "BRL";
  estado: InvoiceEstado;
  dueDate: Date;
  paidAt?: Date;
  createdAt: Date;
  updatedAt: Date;
}

// Para crear: sin campos autogenerados
type InvoiceCreate = Omit<Invoice, "id" | "estado" | "paidAt" | "createdAt" | "updatedAt">;

// Para actualizar: todo opcional
type InvoiceUpdate = Partial<Pick<Invoice, "externalId" | "amount" | "dueDate">>;
\end{lstlisting}
\end{solucionbox}

% ─── CHEATSHEET ──────────────────────────────────────────────────────────────
\newpage
\section{Cheatsheet: TypeScript en Una Página}

\begin{multicols}{2}
\textbf{\textcolor{tsblue}{Tipos básicos}}
\begin{lstlisting}[style=typescript,numbers=none]
string, number, boolean
null, undefined, void
any, unknown, never
T[], Array<T>
[T, U]          // tupla
Record<K, V>    // diccionario
\end{lstlisting}

\textbf{\textcolor{tsblue}{Interfaces}}
\begin{lstlisting}[style=typescript,numbers=none]
interface A { x: number; y?: string }
type B = A & { z: boolean }  // interseccion
type C = A | B               // union
type D = Partial<A>          // todos opcionales
type E = Required<A>         // todos requeridos
type F = Pick<A, "x">       // solo algunos campos
type G = Omit<A, "y">       // sin algunos campos
\end{lstlisting}

\textbf{\textcolor{tsblue}{Funciones}}
\begin{lstlisting}[style=typescript,numbers=none]
function f(x: number): string {}
const f = (x: number): string => {};
const f = (x = 0) => x;      // default
const f = (...xs: number[]) => {}; // rest
\end{lstlisting}

\textbf{\textcolor{tsblue}{Async}}
\begin{lstlisting}[style=typescript,numbers=none]
async function f(): Promise<T> {
  const x = await promesa;
}
await Promise.all([p1, p2]);
await Promise.allSettled([p1, p2]);
\end{lstlisting}

\columnbreak

\textbf{\textcolor{tsblue}{Generics}}
\begin{lstlisting}[style=typescript,numbers=none]
function id<T>(x: T): T { return x; }
class Repo<T> { items: T[] = []; }
function fn<T extends {id: string}>(x: T) {}
\end{lstlisting}

\textbf{\textcolor{vuegreen}{Vue Composition API}}
\begin{lstlisting}[style=typescript,numbers=none]
const x = ref<number>(0);  // reactivo
x.value++;                  // acceso con .value
const obj = reactive({});   // objeto reactivo
const c = computed(() => x.value * 2);
watch(x, (nuevo, viejo) => {});
onMounted(async () => { ... });
\end{lstlisting}

\textbf{\textcolor{vuegreen}{Template Vue}}
\begin{lstlisting}[style=typescript,numbers=none]
{{ variable }}
:prop="valor"    // prop dinamica
@click="fn"      // evento
v-if="cond"
v-for="x in xs" :key="x.id"
v-model="ref"    // two-way binding
\end{lstlisting}

\textbf{\textcolor{warnambel}{Patrones útiles}}
\begin{lstlisting}[style=typescript,numbers=none]
x ?? "default"        // nullish coalescing
x?.propiedad          // optional chaining
x!.propiedad          // non-null assertion
x as Tipo             // type assertion
typeof x === "string" // type guard
\end{lstlisting}
\end{multicols}

\end{document}
